\section{Model Agnostic：没有模型宗教，但要有评测纪律}

从行业案例继续向上，讲者提出 model agnostic：不因 open/closed、small/large 或某个 vendor 建立身份认同，而按生命、效率和生产率目标选择模型。这个立场不是“模型都一样”，而是要求建立更严格的 task contract、evaluation、routing、fallback 和 version governance。

\subsection{Portfolio selection 而非一次性 winner}

本节把“no model religion”落实成持续选择流程，而不是一句供应商中立口号。读图时先从 task contract 出发，再比较候选、证据和 routing policy；这样即使模型能力和价格快速变化，系统仍能解释为什么某类请求使用某个模型、何时 fallback、怎样回滚，以及哪些治理约束不可被性能提升覆盖。

\lecturefigure{10-model-agnostic.png}{Model-agnostic 策略以任务、候选组合、证据和 routing policy 为主线。}{本地字幕 20:20--23:40；概念重绘。}

\begin{knowledgebox}{读图：四步避免 model religion}
先定义任务质量、latency、privacy、language 和 cost；再建立 open/closed、small/large、general/specialist 候选；随后运行 offline/online eval 与 red-team；最后把选择写成可审计 routing policy。新模型进入时重新比较，而不是自动替换所有路径。
\end{knowledgebox}

\begin{importantbox}{模型选择是约束优化}
可写为
\[
m^*=\arg\max_m Q(m)\quad
\text{s.t.}\quad L(m)\leq L_{max},\ C(m)\leq C_{max},\ G(m)=1,
\]
其中 $Q$ 是任务质量，$L$ 是 latency，$C$ 是成本，$G$ 表示治理、数据和部署约束是否满足。若任务分布不同，最优解可以是 routing 或 ensemble，而不是单个模型。
\end{importantbox}

\begin{warningbox}{Model agnostic 不等于供应链无关}
模型可以替换，tokenizer、tool schema、safety policy、observability、fine-tuning data 和 evaluation history 却未必可移植。真正的 portability 需要 versioned interface、prompt contract、data export、fallback 和 rollback，而不是在 UI 中增加一个下拉菜单。
\end{warningbox}

\teachervoice{讲者用“no technology religion”延伸出“no model religion”。这句话的工程版本是：对架构保持可证伪，对供应商保持可替换，对每次迁移保留可比较证据。}

\subsection{本章小结}

Model agnostic 是 evaluation-driven portfolio，而非“任何模型都行”。任务契约、约束、版本和 routing 共同决定选择；越频繁更换模型，越需要稳定的系统接口和回归测试。

